\begin{abstract}
\noindent
After a disaster, alerts raised in the field must reach the rescue center before their deadlines even though parts of the terrestrial network have failed. Existing unmanned aerial vehicle (UAV) relay designs optimize rates or the number of served devices and treat the ground network as an ideal sink, ignoring both per-alert deadlines and the limited backhaul of the surviving relays. This paper jointly optimizes the trajectory of a single UAV relay, the path of every alert through a capacity-limited ground network, and the bandwidth allocation, so as to maximize the fraction of alerts delivered on time under a probabilistic line-of-sight (PrLoS) channel. We prove that the problem is NP-hard even for a fixed trajectory and derive a linear programming (LP) upper bound that holds for every channel realization. We then propose a block search over paths, bandwidth weights, and trajectories whose blocks are local searches around the current plan: paths are rerouted around the links that the plan congests, and trajectory segments are shifted toward saturated UAV links within the speed limit. Every neighbor is scored by simulating the plan on sampled channel realizations, and a variant additionally repairs the schedule of alerts at risk of missing their deadlines. On \ResultTwoScenarioCount{} evaluation scenarios synthesized from Topology Zoo and SNDlib, the block search delivers \ResultTwoTimelyVzeroBtwo{} of the alerts on time, against \ResultTwoTimelyVzeroBone{} for a UAV on a fixed trajectory and \ResultTwoTimelyVzeroBzero{} without a UAV. Its advantage over the fixed trajectory holds on all ten topologies, with a Holm-corrected $p$-value of \ResultTwoHolmVzeroBtwoBone, and grows when the backhaul is narrow. Ablations attribute the gain to selecting paths and trajectory jointly and to designing on sampled realizations rather than expected rates, whereas schedule repair changes the timely-delivery ratio by only \ResultTwoDiffVzeroPBtwo{} and a leximin fairness objective lowers it by \ResultTwoDiffVzeroBtwoBthree. The block search also exceeds an adaptation of the method of Tran et al. by \ResultExtDiffVzeroBtwoTran{} on the default scenarios and by \ResultExtDiffVzeroBhfivezeroBtwoTran{} under narrow backhaul, although neither difference survives Holm correction, and the ranking of all methods is unchanged in a case study with \ResultExtTargetCount{} targets extracted from RescueNet imagery.

\vspace{0.5em}
\noindent
\textbf{Keywords:} UAV relay; trajectory optimization; relay selection; bandwidth allocation; deadline-aware delivery; linear programming bound; local search; paired statistical testing; RescueNet.
\end{abstract}
